\begin{proof}
\begin{enumerate}[label={$\langle1\rangle$\arabic*.},leftmargin=*]
\item 
\noindent\textbf{Step 1.} Apply \cref{cdp:st:A.1.1} to the whole reduced
Cartier fibre \(W_c\), not separately to its components.  Conditions \eqref{wbII:eq:W:5.31}
give
\[
 H^0(W_c,\Omega_X^1|_{W_c}\otimes A_c)
 =\ker\delta_c=0.
\tag{W.59}\hypertarget{tag-03HV}{}\label{wbII:eq:W:5.32}
\]
The map \(f\) is flat: a local parameter at \(c\) pulls back to a
non-zero-divisor in every regular local ring of the connected manifold \(X\),
and torsion-freeness over the base DVR is flatness.  Moreover, the normal
line of the full fibre is the pullback of the one-dimensional vector space
\(T_cC\), hence is trivial.  Adjunction therefore identifies
\(\omega_{W_c}\simeq\omega_X|_{W_c}\).  With \(E=T_X\otimes L\), absolute
Cartier-fibre duality now gives
\[
 H^2(W_c,E|_{W_c})^\vee
 \simeq H^0(W_c,\Omega_X^1|_{W_c}\otimes A_c)=0.
\]
The uniform top-degree base-change isomorphism of
\cref{wbII:cor:surface-top-base-change} therefore gives
\[
 (R^2f_*E)_c\otimes_{\mathcal O_{C,c}}\mathbb C(c)=0
\]
for every \(c\).  Nakayama's lemma kills every stalk of the coherent sheaf
\(R^2f_*E\), proving \eqref{wbII:eq:W:5.29}.  The remaining conclusions follow from
\cref{wbII:cor:corrected-cdp-relative}.

\end{enumerate}
\end{proof}
